
\documentclass[11pt]{article}

\usepackage[margin=1in]{geometry}
\usepackage[utf8]{inputenc}
\usepackage[T1]{fontenc}
\usepackage{lmodern}
\usepackage{amsmath,amssymb,mathtools}
\usepackage{microtype}
\usepackage{hyperref}
\usepackage{enumitem}
\usepackage{fvextra}

% --- Make common Lean unicode compile under pdfLaTeX ---
\DeclareUnicodeCharacter{00AC}{\ensuremath{\neg}}         % ¬
\DeclareUnicodeCharacter{00B7}{\ensuremath{\cdot}}        % ·
\DeclareUnicodeCharacter{2016}{\ensuremath{\Vert}}        % ‖
\DeclareUnicodeCharacter{2080}{\textsubscript{0}}         % ₀
\DeclareUnicodeCharacter{2081}{\textsubscript{1}}         % ₁
\DeclareUnicodeCharacter{2082}{\textsubscript{2}}         % ₂
\DeclareUnicodeCharacter{2083}{\textsubscript{3}}         % ₃
\DeclareUnicodeCharacter{2084}{\textsubscript{4}}         % ₄
\DeclareUnicodeCharacter{2085}{\textsubscript{5}}         % ₅
\DeclareUnicodeCharacter{2086}{\textsubscript{6}}         % ₆
\DeclareUnicodeCharacter{2087}{\textsubscript{7}}         % ₇
\DeclareUnicodeCharacter{2088}{\textsubscript{8}}         % ₈
\DeclareUnicodeCharacter{2089}{\textsubscript{9}}         % ₉
\DeclareUnicodeCharacter{2102}{\ensuremath{\mathbb{C}}}   % ℂ
\DeclareUnicodeCharacter{211D}{\ensuremath{\mathbb{R}}}   % ℝ
\DeclareUnicodeCharacter{2190}{\ensuremath{\leftarrow}}   % ←
\DeclareUnicodeCharacter{2192}{\ensuremath{\to}}          % →
\DeclareUnicodeCharacter{2200}{\ensuremath{\forall}}      % ∀
\DeclareUnicodeCharacter{2203}{\ensuremath{\exists}}      % ∃
\DeclareUnicodeCharacter{2208}{\ensuremath{\in}}          % ∈
\DeclareUnicodeCharacter{2227}{\ensuremath{\wedge}}       % ∧
\DeclareUnicodeCharacter{2228}{\ensuremath{\vee}}         % ∨
\DeclareUnicodeCharacter{2260}{\ensuremath{\neq}}         % ≠
\DeclareUnicodeCharacter{2264}{\ensuremath{\le}}          % ≤
\DeclareUnicodeCharacter{2265}{\ensuremath{\ge}}          % ≥
\DeclareUnicodeCharacter{27E8}{\ensuremath{\langle}}      % ⟨
\DeclareUnicodeCharacter{27E9}{\ensuremath{\rangle}}      % ⟩

% Verbatim environment for Lean code (line-breaking enabled)
\DefineVerbatimEnvironment{LeanCode}{Verbatim}{
  fontsize=\small,
  breaklines=true,
  breakanywhere=true,
  commandchars=\\\{\},
  showspaces=false,
  showtabs=false
}

\setlist[itemize]{topsep=2pt, itemsep=2pt}
\setlist[enumerate]{topsep=2pt, itemsep=2pt}

\title{Detailed Notes on \texttt{D3Sec17.lean}\\(No real solution to a 19-equation quadratic system)}
\author{}
\date{}

\begin{document}
\maketitle

\section*{-- general overview and logic: what is the problem, how does the proof go (and check if the proof is logically correct)}

\subsection*{What is being proved?}

The file proves a final theorem:

\begin{quote}
\textbf{No real solution exists} to a specific system of \textbf{19 quadratic equations} in \textbf{28 real variables}.
\end{quote}

Concretely, the theorem at the end is (in Lean):

\begin{LeanCode}
theorem NoRealSolution : ¬ ∃ v : Fin 28 → ℝ,
  let c := c_from_v v
  E1 c ∧ E2 c ∧ ... ∧ E19 c
\end{LeanCode}

This means: there is no vector $v\in\mathbb{R}^{28}$ such that (after mapping $v$ into $14$ complex numbers $c_0,\dots,c_{13}$) all equations $E1,\dots,E19$ hold.

\subsection*{How are the variables organized?}

Instead of working directly with $28$ real variables, the file packages them into $14$ complex variables:
\[
c_0,\dots,c_{13}\in\mathbb{C}.
\]
Each $c_k$ has real and imaginary parts coming from two real coordinates.  The key derived quantity is
\[
p_k := \|c_k\|^2 = |c_k|^2 \in \mathbb{R}_{\ge 0},
\]
implemented as \texttt{Complex.normSq (c k)} in Lean.

So the system is quadratic in the reals, but becomes a structured system involving:
\begin{itemize}
\item linear equations in the nonnegative reals $p_k$, and
\item bilinear equations involving products like $\overline{c_i}c_j$.
\end{itemize}

\subsection*{What do the 19 equations look like?}

They split naturally into blocks.

\paragraph{(i) ``Diagonal'' / norm-squared constraints (E1--E8).}
These are linear equalities among the $p_k = |c_k|^2$. For example:
\[
\sum_{k=0}^{13} p_k = 1,\qquad
p_7+p_8+\cdots+p_{13}=\frac12,
\]
and six other similar subset sums.

\paragraph{(ii) A ``Y1'' imaginary-part sum (E9).}
Equation E9 says a sum of imaginary parts vanishes:
\[
\operatorname{Im}(\overline{c_0}c_7)+\cdots+\operatorname{Im}(\overline{c_6}c_{13})=0.
\]
Later the proof splits this into two pieces $U$ (indices $0$--$3$) and $T$ (indices $4$--$6$).

\paragraph{(iii) Local orthogonality-type constraints (E10--E15).}
These constrain real and imaginary parts of
\[
\overline{c_5}c_{13} + \overline{c_6}c_{12}, \qquad
\overline{c_4}c_{13} + \overline{c_6}c_{11},
\]
and force corresponding multiplicative relations among norms.

\paragraph{(iv) Real-part constraints (E16--E19).}
Together these force certain real parts to vanish, e.g.
\[
\operatorname{Re}(\overline{c_{12}}c_{13})=0,\qquad
\operatorname{Re}(\overline{c_{5}}c_{6})=0,
\]
and similarly for $(c_{11},c_{13})$ and $(c_4,c_6)$.

\subsection*{How does the proof work? (main logical chain)}

The proof is a sequence of lemmas \texttt{L1}--\texttt{L11} plus a final gauge argument in \texttt{NoRealSolution}.  The key idea is:

\begin{enumerate}
\item Use the diagonal equations (E1--E8) to solve for many combinations of the $p_k$ (Lemma L1).
\item Use the bilinear blocks (E10--E12) and (E13--E15) to derive two \emph{multiplicative identities} among the $p_k$ (L2, L3).
\item Combine (1) and (2) to get \emph{exact proportionalities}:
  \[
  p_{11}=p_{12}=\frac35 p_{13},\qquad p_4=p_5=\frac35 p_6 \quad \text{(Lemma L4)}.
  \]
\item From those proportionalities and the diagonal relations in L1, derive \emph{uniform lower bounds}:
  \[
  p_{13}\ge \frac{15}{176},\qquad p_6\ge \frac{15}{176} \quad \text{(Lemma L5)}.
  \]
\item Perform a \textbf{global phase normalization} (a gauge transformation): multiply all $c_k$ by a unit complex number $g$ so that $c_{13}$ becomes a positive real. This does not change any $p_k$ nor any inner-products $\overline{c_i}c_j$ (Lemma L7).
\item Under the gauge $c_{13}\in\mathbb{R}_{>0}$, the real-part equations force certain variables to be purely real/purely imaginary (Lemma L8).
\item Define
  \[
  T := \operatorname{Im}(\overline{c_4}c_{11})+\operatorname{Im}(\overline{c_5}c_{12})+\operatorname{Im}(\overline{c_6}c_{13})
  \]
  and show
  \[
  |T| = \frac{11}{5}\sqrt{p_6p_{13}} \ \ge\ \frac{3}{16} \quad \text{(Lemma L9)}.
  \]
\item Define
  \[
  U := \operatorname{Im}(\overline{c_0}c_{7})+\operatorname{Im}(\overline{c_1}c_{8})+\operatorname{Im}(\overline{c_2}c_{9})+\operatorname{Im}(\overline{c_3}c_{10})
  \]
  and show
  \[
  |U|\le \frac{5}{32} \quad \text{(Lemma L10)}.
  \]
\item Equation E9 is exactly $U+T=0$, hence $|T|=|U|$ (Lemma L11).
\item Contradiction:
  \[
  |T|\ge \frac{3}{16}=\frac{6}{32} > \frac{5}{32}\ge |U|,\quad \text{but } |T|=|U|.
  \]
\end{enumerate}

\subsection*{Is the proof logically correct?}

\begin{itemize}
\item \textbf{Mathematical correctness (by inspection):} Each step is standard:
  \begin{itemize}
  \item L1 is linear elimination on the diagonal equations.
  \item L2/L3 use that the block equations force two complex numbers to be negatives of each other and purely real, hence their squared magnitudes are equal.
  \item L4 is algebra combining L1 with L2/L3.
  \item L5 is inequality chasing using nonnegativity of norm squares.
  \item L7 is a standard unit-modulus ``global phase'' gauge.
  \item L8--L11 are structured real/imag decompositions plus the triangle/AM--GM style inequality $| \operatorname{Im}(\overline{a}b) |\le (|a|^2+|b|^2)/2$.
  \end{itemize}
\item \textbf{Lean correctness:} If \texttt{D3Sec17.lean} typechecks in Lean, then the final theorem is certified by Lean's kernel. This file contains no \texttt{sorry} and no \texttt{exact?} placeholders; it relies on automation (\texttt{grind}, \texttt{linarith}, \texttt{nlinarith}) but those tactics produce explicit proof terms when they succeed.
\end{itemize}

\newpage
\section*{-- interpretation of lean code: a few lines of lean code, then explanation; then next few lines of lean code, then explanation; \dots}

\subsection*{Imports, options, and the basic ``norm squared'' function}

\begin{LeanCode}
import Mathlib

noncomputable section

def p (c : Fin 14 → ℂ) (k : Fin 14) : ℝ := Complex.normSq (c k)
\end{LeanCode}

\paragraph{Explanation.}
\texttt{p c k} is the real number $\|c_k\|^2 = |c_k|^2$.
\texttt{Complex.normSq} returns a nonnegative real number and satisfies $\mathrm{normSq}(zw)=\mathrm{normSq}(z)\,\mathrm{normSq}(w)$.

\subsection*{The diagonal equations E1--E8}

\begin{LeanCode}
def E1 (c : Fin 14 → ℂ) : Prop :=
  (Finset.univ.sum (fun i => p c i)) = 1

def E2 (c : Fin 14 → ℂ) : Prop :=
  p c 7 + p c 8 + p c 9 + p c 10 + p c 11 + p c 12 + p c 13 = 1/2

def E3 (c : Fin 14 → ℂ) : Prop :=
  p c 5 + p c 6 + p c 12 + p c 13 = 1/2
\end{LeanCode}

\paragraph{Explanation.}
\begin{itemize}
\item E1 says the total sum of squared magnitudes is $1$:
  \[
  \sum_{k=0}^{13} p_k = 1.
  \]
\item E2 says the ``tail'' sum $p_7+\cdots+p_{13}$ equals $1/2$.
\item E3 (and similarly E4--E8, not shown here) constrain other subset sums to equal $1/2$.
\end{itemize}

Together E1--E8 form a linear system in the nonnegative reals $p_0,\dots,p_{13}$.

\subsection*{The Y1 equation E9 and the local bilinear blocks}

\begin{LeanCode}
def E9 (c : Fin 14 → ℂ) : Prop :=
  (star (c 0) * c 7).im + (star (c 1) * c 8).im +
  (star (c 2) * c 9).im + (star (c 3) * c 10).im +
  (star (c 4) * c 11).im + (star (c 5) * c 12).im +
  (star (c 6) * c 13).im = 0
\end{LeanCode}

\paragraph{Explanation.}
Here \texttt{star} is complex conjugation, so \texttt{star (c i) * c j} is $\overline{c_i}c_j$.
Thus E9 is:
\[
\sum_{i=0}^{6} \operatorname{Im}(\overline{c_i}c_{i+7}) = 0.
\]
Later this is split into $U$ (the first four terms) and $T$ (the last three terms).

\begin{LeanCode}
def E10 (c : Fin 14 → ℂ) : Prop :=
  (star (c 5) * c 13).re + (star (c 6) * c 12).re = 0

def E11 (c : Fin 14 → ℂ) : Prop :=
  (star (c 5) * c 13).im + (star (c 6) * c 12).im = 0

def E12 (c : Fin 14 → ℂ) : Prop :=
  (star (c 5) * c 13).im - (star (c 6) * c 12).im = 0
\end{LeanCode}

\paragraph{Explanation.}
Let $A=\overline{c_5}c_{13}$ and $B=\overline{c_6}c_{12}$.
\begin{itemize}
\item E10 says $\operatorname{Re}(A)+\operatorname{Re}(B)=0$.
\item E11 says $\operatorname{Im}(A)+\operatorname{Im}(B)=0$.
\item E12 says $\operatorname{Im}(A)-\operatorname{Im}(B)=0$.
\end{itemize}
From E11 and E12 we get $\operatorname{Im}(A)=\operatorname{Im}(B)=0$.
Together with E10 this implies $A=-B$ (as real numbers inside $\mathbb{C}$), which is the input to Lemma L2.

\subsection*{Lemma L1: solving the diagonal subsystem}

\begin{LeanCode}
theorem L1 (c : Fin 14 → ℂ)
  (h1 : E1 c) (h2 : E2 c) (h3 : E3 c) (h4 : E4 c)
  (h5 : E5 c) (h6 : E6 c) (h7 : E7 c) (h8 : E8 c) :
  p c 0 + p c 7 = 1/16 ∧
  p c 1 + p c 8 = 1/16 ∧
  p c 2 + p c 9 = 1/16 ∧
  p c 3 + p c 10 = 1/8 ∧
  p c 4 + p c 11 = 3/16 ∧
  p c 5 + p c 12 = 3/16 ∧
  p c 6 + p c 13 = 5/16 ∧
  p c 7 = 1/2 - (p c 8 + p c 9 + p c 10 + p c 11 + p c 12 + p c 13) ∧
  p c 0 = (p c 8 + p c 9 + p c 10 + p c 11 + p c 12 + p c 13) - 7/16 := by
  unfold E1 E2 E3 E4 E5 E6 E7 E8 at *; norm_num [ Fin.sum_univ_succ ] at *;
  grind
\end{LeanCode}

\paragraph{Explanation.}
Lemma L1 is purely linear algebra on the equations E1--E8.
It yields:
\begin{itemize}
\item \emph{Pair-sum identities} such as $p_0+p_7=\frac1{16}$, $p_1+p_8=\frac1{16}$, etc.
\item \emph{Explicit formulas} for $p_7$ and $p_0$ in terms of $(p_8,p_9,p_{10},p_{11},p_{12},p_{13})$.
\end{itemize}
The proof is handled automatically by \texttt{grind} after expanding the definitions.

\subsection*{Lemmas L2 and L3: multiplicative identities from the bilinear blocks}

\begin{LeanCode}
theorem L2 (c : Fin 14 → ℂ)
  (h10 : E10 c) (h11 : E11 c) (h12 : E12 c) :
  p c 5 * p c 13 = p c 6 * p c 12 := by
  unfold E10 E11 E12 at *;
  unfold p;
  simp_all +decide [Complex.normSq];
  grind +ring
\end{LeanCode}

\paragraph{Explanation.}
As discussed above, E10--E12 force $\overline{c_5}c_{13}=-\overline{c_6}c_{12}$ and both sides are real.
Taking squared magnitudes and using multiplicativity of \texttt{normSq} gives:
\[
|c_5|^2\,|c_{13}|^2 = |c_6|^2\,|c_{12}|^2,
\]
i.e.\ $p_5p_{13}=p_6p_{12}$.

\begin{LeanCode}
theorem L3 (c : Fin 14 → ℂ)
  (h13 : E13 c) (h14 : E14 c) (h15 : E15 c) :
  p c 4 * p c 13 = p c 6 * p c 11 := by
  unfold E13 E14 E15 p at *;
  simp_all +decide [Complex.normSq];
  grind
\end{LeanCode}

\paragraph{Explanation.}
This is the analogous statement for the block E13--E15, giving $p_4p_{13}=p_6p_{11}$.

\subsection*{Lemma L4: deducing exact proportionalities among norms}

\begin{LeanCode}
theorem L4 (c : Fin 14 → ℂ)
  (h1 : E1 c) (h2 : E2 c) (h3 : E3 c) (h4 : E4 c)
  (h5 : E5 c) (h6 : E6 c) (h7 : E7 c) (h8 : E8 c)
  (h10 : E10 c) (h11 : E11 c) (h12 : E12 c)
  (h13 : E13 c) (h14 : E14 c) (h15 : E15 c) :
  p c 12 = 3/5 * p c 13 ∧
  p c 11 = 3/5 * p c 13 ∧
  p c 5 = 3/5 * p c 6 ∧
  p c 4 = 3/5 * p c 6 := by
  ...
  grind
\end{LeanCode}

\paragraph{Explanation.}
The goal is to solve the multiplicative identities from L2/L3 together with the pair-sum identities from L1.
Sketch of the algebra (informal):
\begin{itemize}
\item From L1: $p_5+p_{12}=\frac{3}{16}$ and $p_6+p_{13}=\frac{5}{16}$.
\item From L2: $p_5p_{13}=p_6p_{12}$.
\end{itemize}
This system has the solution $p_{12}=\frac35 p_{13}$ and $p_5=\frac35 p_6$ (and similarly with $p_{11},p_4$ using L3).
Lean lets \texttt{grind} perform the necessary algebra.

\subsection*{Lemma L5: uniform lower bounds on $p_{13}$ and $p_6$}

\begin{LeanCode}
theorem L5 (c : Fin 14 → ℂ)
  (h1 : E1 c) (h2 : E2 c) (h3 : E3 c) (h4 : E4 c)
  (h5 : E5 c) (h6 : E6 c) (h7 : E7 c) (h8 : E8 c)
  (h10 : E10 c) (h11 : E11 c) (h12 : E12 c)
  (h13 : E13 c) (h14 : E14 c) (h15 : E15 c) :
  p c 13 ≥ 15/176 ∧ p c 6 ≥ 15/176 := by
  ...
  constructor <;> linarith [...]
\end{LeanCode}

\paragraph{Explanation.}
This lemma is an inequality chase using:
\begin{itemize}
\item L4: $p_{11}=p_{12}=\frac35 p_{13}$,
\item L1: $p_0 = (p_8+p_9+p_{10}+p_{11}+p_{12}+p_{13})-\frac{7}{16}$,
\item nonnegativity $p_k\ge 0$ for all $k$.
\end{itemize}

Key bound for $p_{13}$:
\[
0\le p_0 \implies p_8+p_9+p_{10}+p_{11}+p_{12}+p_{13}\ge \frac{7}{16}.
\]
Substitute $p_{11}=p_{12}=\frac35 p_{13}$ to get
\[
p_8+p_9+p_{10}+\frac{11}{5}p_{13}\ge \frac{7}{16}.
\]
Also L1 gives $p_1+p_8=\frac1{16}$, $p_2+p_9=\frac1{16}$, $p_3+p_{10}=\frac18$, so with $p_1,p_2,p_3\ge 0$:
\[
p_8\le \frac{1}{16},\quad p_9\le \frac{1}{16},\quad p_{10}\le \frac{1}{8}.
\]
Hence $p_8+p_9+p_{10}\le \frac14$, so
\[
\frac{11}{5}p_{13}\ge \frac{7}{16}-\frac14=\frac{3}{16}
\quad\Longrightarrow\quad
p_{13}\ge \frac{3}{16}\cdot\frac{5}{11}=\frac{15}{176}.
\]
A similar argument gives an upper bound $p_{13}\le \frac{5}{22}$ from $p_7\ge 0$, and since L1 gives $p_6=\frac{5}{16}-p_{13}$, this yields
\[
p_6\ge \frac{15}{176}.
\]
Lean discharges these inequalities with \texttt{linarith}.

\subsection*{Lemma L6: consequences of E16--E19}

\begin{LeanCode}
theorem L6 (c : Fin 14 → ℂ)
  (h16 : E16 c) (h17 : E17 c) (h18 : E18 c) (h19 : E19 c) :
  (star (c 12) * c 13).re = 0 ∧
  (star (c 5) * c 6).re = 0 ∧
  (star (c 11) * c 13).re = 0 ∧
  (star (c 4) * c 6).re = 0 := by
  unfold E16 E17 E18 E19 at *;
  grind +ring
\end{LeanCode}

\paragraph{Explanation.}
For example, E16 and E17 are:
\[
\operatorname{Re}(\overline{c_{12}}c_{13})+\operatorname{Re}(\overline{c_5}c_6)=0,\qquad
\operatorname{Re}(\overline{c_{12}}c_{13})-\operatorname{Re}(\overline{c_5}c_6)=0.
\]
Adding and subtracting gives each term is $0$.  Similarly for E18/E19.
Lemma L6 packages these consequences.

\subsection*{Lemma L7: global phase normalization (gauge)}

\begin{LeanCode}
theorem L7 (c : Fin 14 → ℂ)
  (h1 : E1 c) ... (h15 : E15 c) :
  let g := star (c 13) / (‖c 13‖ : ℂ)
  let c' := fun k => g * c k
  c' 13 = (‖c 13‖ : ℂ) ∧
  (∀ i j, star (c' i) * c' j = star (c i) * c j) ∧
  (∀ k, p c' k = p c k) := by
  by_cases h : c 13 = 0 <;> simp_all +decide [ div_eq_inv_mul ];
  · have := L5 c h1 ... h15; ...
  · field_simp [h]; ...
\end{LeanCode}

\paragraph{Explanation.}
The gauge choice is
\[
g:=\frac{\overline{c_{13}}}{\|c_{13}\|}.
\]
Then $g$ has norm $1$, and
\[
c'_{13} = g c_{13} = \|c_{13}\|\in\mathbb{R}_{\ge 0}.
\]
Moreover,
\[
\overline{(g c_i)}(g c_j)=\overline{c_i}c_j
\]
because $\overline{g}g=1$, so all bilinear expressions $\overline{c_i}c_j$ are invariant under this global rotation.  Finally, since $|g|=1$, the norms satisfy $|g c_k|^2=|c_k|^2$, i.e.\ $p$ is unchanged.

The proof splits into cases: if $c_{13}=0$ then $p_{13}=0$, contradicting L5's lower bound; so the nonzero case holds and the algebra goes through.

\subsection*{Lemma L8: under the gauge $c_{13}\in\mathbb{R}_{>0}$, many variables become purely real/imaginary}

\begin{LeanCode}
theorem L8 (c : Fin 14 → ℂ)
  (h1 : E1 c) ... (h19 : E19 c)
  (h_gauge_im : (c 13).im = 0) (h_gauge_pos : 0 < (c 13).re) :
  (c 4).im = 0 ∧ (c 5).im = 0 ∧
  (c 11).re = 0 ∧ (c 12).re = 0 ∧
  (c 6).re = 0 ∧
  c 4 ≠ 0 ∧ c 5 ≠ 0 ∧ (c 6).im ≠ 0 := by
  ...
\end{LeanCode}

\paragraph{Explanation.}
Under the gauge assumptions, $c_{13}$ is a positive real number. Then:
\begin{itemize}
\item From E14 and E15, one extracts $\operatorname{Im}(\overline{c_4}c_{13})=0$, and since $c_{13}\neq 0$ and is real, this forces $\operatorname{Im}(c_4)=0$. Hence $c_4$ is real.
\item From L6, $\operatorname{Re}(\overline{c_{11}}c_{13})=0$ and $\operatorname{Re}(\overline{c_{12}}c_{13})=0$, which force $\operatorname{Re}(c_{11})=\operatorname{Re}(c_{12})=0$ (so $c_{11},c_{12}$ are purely imaginary).
\item Another real-part equation from L6 gives $\operatorname{Re}(\overline{c_4}c_6)=0$, and since $c_4$ is real and nonzero, this forces $\operatorname{Re}(c_6)=0$ (so $c_6$ is purely imaginary).
\item The nonzero conditions $c_4\neq 0$, $c_5\neq 0$, and $(c_6).im\neq 0$ are proved using L4 and L5: if one of these vanished, it would force a corresponding norm-square to be $0$, contradicting the proportionalities and the lower bounds.
\end{itemize}

\subsection*{Lemma L9: define $T$ and compute $|T|$ and a lower bound}

\begin{LeanCode}
theorem L9 (c : Fin 14 → ℂ)
  (h1 : E1 c) ... (h19 : E19 c)
  (h_gauge_im : (c 13).im = 0) (h_gauge_pos : 0 < (c 13).re) :
  let T := (star (c 4) * c 11).im + (star (c 5) * c 12).im + (star (c 6) * c 13).im
  abs T = 11/5 * Real.sqrt (p c 6 * p c 13) ∧
  abs T ≥ 3/16 := by
  ...
\end{LeanCode}

\paragraph{Explanation.}
Under the gauge and using L8, write:
\[
c_4=y_4\in\mathbb{R},\quad c_5=y_5\in\mathbb{R},\quad
c_{11}= i y_{11},\quad c_{12}= i y_{12},\quad c_6= i y_6,\quad c_{13}=r_{13}\in\mathbb{R}_{>0}.
\]
Then:
\[
\operatorname{Im}(\overline{c_4}c_{11}) = \operatorname{Im}(y_4 \cdot i y_{11}) = y_4 y_{11},
\]
\[
\operatorname{Im}(\overline{c_5}c_{12}) = y_5 y_{12},\qquad
\operatorname{Im}(\overline{c_6}c_{13}) = \operatorname{Im}((-i y_6) r_{13}) = -y_6 r_{13}.
\]
So $T = y_4 y_{11} + y_5 y_{12} - y_6 r_{13}$.

The remaining equations force $y_{11}$ and $y_{12}$ to be affine multiples of $y_6 r_{13}$, and the file shows:
\[
T = \pm \frac{11}{5}\, y_6 r_{13}.
\]
Thus
\[
|T| = \frac{11}{5}\,|y_6|\,r_{13}.
\]
But $p_6=|c_6|^2=y_6^2$ and $p_{13}=|c_{13}|^2=r_{13}^2$, so
\[
|y_6|\,r_{13} = \sqrt{y_6^2}\,\sqrt{r_{13}^2} = \sqrt{p_6p_{13}},
\]
giving the formula $|T|=\frac{11}{5}\sqrt{p_6p_{13}}$.

Finally, Lemma L5 gives $p_6,p_{13}\ge \frac{15}{176}$, hence
\[
|T|=\frac{11}{5}\sqrt{p_6p_{13}}
\ge \frac{11}{5}\sqrt{\Bigl(\frac{15}{176}\Bigr)^2}
= \frac{11}{5}\cdot\frac{15}{176}
= \frac{3}{16}.
\]

\subsection*{Lemma L10: define $U$ and bound it from above}

\begin{LeanCode}
theorem L10 (c : Fin 14 → ℂ)
  (h1 : E1 c) ... (h8 : E8 c) :
  let U := (star (c 0) * c 7).im + (star (c 1) * c 8).im
         + (star (c 2) * c 9).im + (star (c 3) * c 10).im
  abs U ≤ 5/32 := by
  ...
\end{LeanCode}

\paragraph{Explanation.}
For any complex numbers $a,b$:
\[
|\operatorname{Im}(\overline{a}b)| \le |\overline{a}b| = |a||b|
\le \frac{|a|^2+|b|^2}{2}.
\]
In the file, this is applied term-by-term to $U$, yielding:
\[
|U|
\le \sum_{i=0}^{3} \frac{p_i+p_{i+7}}{2}.
\]
Lemma L1 provides the exact pair sums:
\[
p_0+p_7=\frac{1}{16},\quad p_1+p_8=\frac{1}{16},\quad p_2+p_9=\frac{1}{16},\quad p_3+p_{10}=\frac{1}{8}.
\]
Therefore
\[
|U| \le \frac{1}{2}\Bigl(\frac{1}{16}+\frac{1}{16}+\frac{1}{16}+\frac{1}{8}\Bigr)
= \frac{1}{2}\Bigl(\frac{3}{16}+\frac{2}{16}\Bigr)
= \frac{5}{32}.
\]

\subsection*{Lemma L11: E9 ties $U$ and $T$}

\begin{LeanCode}
theorem L11 (c : Fin 14 → ℂ)
  (h9 : E9 c) :
  let T := (star (c 4) * c 11).im + (star (c 5) * c 12).im + (star (c 6) * c 13).im
  let U := (star (c 0) * c 7).im + (star (c 1) * c 8).im + (star (c 2) * c 9).im + (star (c 3) * c 10).im
  U + T = 0 ∧ abs T = abs U := by
  unfold E9 at h9;
  exact ⟨ by linarith, by rw [ ← abs_neg ] ; congr; linarith ⟩
\end{LeanCode}

\paragraph{Explanation.}
E9 is exactly $U+T=0$, so $T=-U$. Hence $|T|=|U|$.
This lemma extracts both facts directly.

\subsection*{Mapping 28 reals to 14 complexes and the final contradiction}

\begin{LeanCode}
def c_from_v (v : Fin 28 → ℝ) : Fin 14 → ℂ :=
  fun k => ⟨v ⟨2 * k.val, by fin_cases k <;> decide⟩,
            v ⟨2 * k.val + 1, by decide +revert⟩⟩
\end{LeanCode}

\paragraph{Explanation.}
This is the packaging map:
\[
c_k = v_{2k} + i\,v_{2k+1}.
\]
So a real vector $v\in\mathbb{R}^{28}$ is identified with $14$ complex numbers.

\begin{LeanCode}
theorem NoRealSolution : ¬ ∃ v : Fin 28 → ℝ,
  let c := c_from_v v
  E1 c ∧ E2 c ∧ E3 c ∧ ... ∧ E19 c := by
  by_contra h
  obtain ⟨v, hc⟩ := h
  set c := c_from_v v
  ...
  have := L9 ( fun k => g * c k ) ... hg₂ hg₃
  have := L10 ( fun k => g * c k ) ...
  have := L11 ( fun k => g * c k ) hc₉
  linarith
\end{LeanCode}

\paragraph{Explanation.}
The final theorem proceeds by contradiction:
\begin{enumerate}
\item Assume there is a real solution $v$; set $c=c\_from\_v(v)$ and collect the hypotheses $E1(c),\dots,E19(c)$.
\item Choose a gauge $g=\overline{c_{13}}/\|c_{13}\|$ so that $(g c_{13})$ is a positive real.
\item Prove that multiplying all variables by $g$ preserves all equations (because $\overline{(g c_i)}(g c_j)=\overline{c_i}c_j$ and $p(gc_k)=p(c_k)$ when $\|g\|=1$).
\item Apply Lemmas L9, L10, L11 to the gauged tuple $c'=\{g c_k\}$:
  \[
  |T|\ge \frac{3}{16},\qquad |U|\le \frac{5}{32},\qquad |T|=|U|.
  \]
\item The numeric contradiction is discharged by \texttt{linarith}.
\end{enumerate}

This completes the proof that no real solution exists.

\end{document}
